%======================
% Appendix
%======================

\section*{A.1 \quad Self-Collected Dataset Protocol}

The SPARK self-collected dataset is recorded at KEC campus using the SPARK wearable node (ESP32 + MPU6050 at 100\,Hz). All four team members plus recruited volunteers act as subjects. The protocol defines four fall types and five ADL types as listed below.

\textbf{Fall types (positive class):}
\begin{itemize}
    \item Forward trip fall --- simulated stumble onto a crash mat
    \item Lateral fall (left) --- sideways fall onto crash mat
    \item Lateral fall (right) --- sideways fall onto crash mat
    \item Sitting-to-floor fall --- controlled seat loss from a chair
\end{itemize}

\textbf{ADL types (negative class):}
\begin{itemize}
    \item Walking (normal pace, indoor)
    \item Running (moderate pace, corridor)
    \item Stair climbing (ascending and descending, KEC stairwell)
    \item Sit-to-stand transition (chair, rapid)
    \item Bending to pick up an object
\end{itemize}

Each fall is repeated a minimum of 5 times per subject. ADL sequences are recorded in 2-minute continuous blocks. The node is worn on the dominant wrist during all recordings. Raw CSV output is labelled using a post-hoc annotation script keyed to the Layer~1 flag timestamp.

\section*{A.2 \quad ESP32 Hardware Connections}

\begin{table}[H]
    \centering
    \caption*{Wearable Node Pin Connections}
    \begin{tabular}{|l|l|l|}
        \hline
        \textbf{Component} & \textbf{ESP32 Pin} & \textbf{Function} \\
        \hline
        MPU6050 SDA & GPIO 21 & I$^2$C Data \\
        MPU6050 SCL & GPIO 22 & I$^2$C Clock \\
        MPU6050 VCC & 3.3\,V rail & Power \\
        MPU6050 GND & GND & Ground \\
        INA219 SDA & GPIO 21 & I$^2$C Data (address 0x41) \\
        INA219 SCL & GPIO 22 & I$^2$C Clock \\
        INA219 VIN+ & After LDO output & Current sense (series) \\
        TP4056 OUT+ & LDO AMS1117 IN & Battery regulated output \\
        AMS1117-3.3 OUT & ESP32 VIN & 3.3\,V regulated supply \\
        \hline
    \end{tabular}
\end{table}

\section*{A.3 \quad PostgreSQL Schema}

\begin{verbatim}
CREATE TABLE fall_events (
    event_id           UUID PRIMARY KEY DEFAULT gen_random_uuid(),
    user_id            VARCHAR(64) NOT NULL,
    timestamp          TIMESTAMPTZ NOT NULL,
    layer1_flag        BOOLEAN NOT NULL,
    layer2_confidence  FLOAT NOT NULL,
    fall_confirmed     BOOLEAN NOT NULL,
    ax_peak            FLOAT,
    ay_peak            FLOAT,
    az_peak            FLOAT,
    severity_score     INT CHECK (severity_score BETWEEN 1 AND 5),
    shap_top_feature   VARCHAR(32),
    shap_json          JSONB,
    alert_sent         BOOLEAN DEFAULT FALSE,
    false_positive_flagged BOOLEAN DEFAULT FALSE
);
\end{verbatim}
